\documentclass[11pt,a4paper]{article}

\usepackage[utf8]{inputenc}
\usepackage[T1]{fontenc}
\usepackage{lmodern}
\usepackage{microtype}
\usepackage{geometry}
\usepackage{amsmath,amssymb,mathtools}
\usepackage{booktabs}
\usepackage{graphicx}
\usepackage{xcolor}
\usepackage{hyperref}
\usepackage{enumitem}

\geometry{margin=2.25cm}
\hypersetup{
  colorlinks=true,
  linkcolor=blue!55!black,
  citecolor=blue!55!black,
  urlcolor=blue!55!black,
  pdftitle={Phase-space CAD for cold-atom quantum sensors},
  pdfauthor={FondamentalWork}
}

\setlength{\parindent}{0pt}
\setlength{\parskip}{0.55em}
\setlist[itemize]{leftmargin=1.35em,topsep=0.25em,itemsep=0.15em}

\newcommand{\kb}{k_{\rm B}}
\newcommand{\lth}{\lambda_{\rm th}}

\title{\textbf{Phase-space CAD for cold-atom quantum sensors}\\
\large A conservative design-budget framework from $D^3v^3\simeq h^3/m^3$}
\author{FondamentalWork}
\date{22 May 2026}

\begin{document}
\maketitle

\begin{abstract}
We turn the position--velocity phase-space cell relation
$D^3v^3\simeq h^3/m^3$ into a practical first-order design framework for
cold-atom quantum sensors. The proposal is not a new propulsion mechanism and
does not rely on experimentally inaccessible GUP corrections. Its useful layer
is a conservative phase-space budget linking atomic temperature, density,
velocity dispersion, cloud expansion, aperture survival, detected atom number
and atom-interferometer shot noise. The resulting software concept,
\texttt{PhaseSpace CAD}, can reject physically weak source architectures before
full optical, vibration and navigation simulations are performed.
\end{abstract}

\section{Purpose}

The most useful industrial application found for the phase-space formulas is
not chemical propulsion. Hot rocket gases have extremely small
$n\lambda_{\rm th}^3$ and are deeply classical. The strong application is
quantum metrology: inertial sensors, gravimeters, gyroscopes and gravity
gradiometers based on cold atoms.

The key design question is simple: can a candidate atom source reach its target
sensitivity within its allowed temperature, density, atom number, package size
and interrogation time? If the answer is negative at the phase-space-budget
level, the detailed instrument cannot recover the lost performance.

\section{Phase-space budget}

For a thermal atomic cloud,
\begin{equation}
\lth=\frac{h}{\sqrt{2\pi m\kb T}},
\qquad
\eta=n\lth^3,
\end{equation}
where $\eta$ is the thermal phase-space density. The one-dimensional velocity
spread is
\begin{equation}
\sigma_v=\sqrt{\frac{\kb T}{m}}.
\end{equation}
For ballistic free expansion,
\begin{equation}
R(t)=\sqrt{R_0^2+\sigma_v^2t^2}.
\end{equation}
These equations provide the source budget. A colder source has larger
$\eta$, smaller $\sigma_v$, slower expansion and better aperture survival.

\section{Interferometer sensitivity}

For a light-pulse atom interferometer, acceleration is encoded as
\begin{equation}
\Delta\phi=k_{\rm eff}aT_i^2.
\end{equation}
The ideal shot-noise acceleration floor is approximated by
\begin{equation}
\delta a\simeq
\frac{1}{k_{\rm eff}T_i^2\sqrt{N_{\rm det}}}.
\end{equation}
For a gyroscope with transverse velocity $v$, a Sagnac-like proxy is
\begin{equation}
\delta\Omega\simeq
\frac{1}{2k_{\rm eff}vT_i^2\sqrt{N_{\rm det}}}.
\end{equation}
Longer interrogation improves the ideal phase response, but it also increases
cloud expansion. The detected atom number therefore depends on the same
temperature that sets velocity dispersion.

\section{Implemented CAD proxy}

The script \texttt{scripts/phase\_space\_sensor\_cad.py} scans temperature and
interrogation time for four use cases:
\begin{itemize}
  \item Rb87 inertial accelerometer for GPS-denied navigation;
  \item Rb87 portable field gravimeter;
  \item Cs133 compact atomic gyroscope;
  \item Sr88 source for space gravity gradiometry.
\end{itemize}

Each design receives a bounded utility score combining aperture survival,
shot-noise target, phase-space density, package compactness and a practical
cooling floor. The cooling term prevents the optimizer from blindly selecting
nanokelvin temperatures when a deployable microkelvin or sub-microkelvin source
would be the more realistic engineering choice.

\begin{table}[h]
\centering
\small
\begin{tabular}{llllrrr}
\toprule
Scenario & Atom & $T$ [K] & $T_i$ [s] & $\eta$ & radius [mm] & score\\
\midrule
GPS-denied accelerometer & Rb87 & $1.0\times10^{-7}$ & 0.05 & $2.08\times10^{-1}$ & 0.81 & 0.974\\
Field gravimeter & Rb87 & $5.0\times10^{-8}$ & 0.25 & $2.93\times10^{-1}$ & 1.32 & 0.980\\
Compact gyroscope & Cs133 & $1.0\times10^{-7}$ & 0.10 & $8.79\times10^{-2}$ & 0.74 & 0.960\\
Space gradiometer source & Sr88 & $1.0\times10^{-8}$ & 1.00 & $1.29$ & 1.79 & 1.000\\
\bottomrule
\end{tabular}
\caption{Best design points from the first-order phase-space CAD proxy. These
are feasibility-screening points, not final instrument specifications.}
\end{table}

\begin{figure}[h]
\centering
\includegraphics[width=0.96\linewidth]{reports/figures/phase_space_sensor_cad.png}
\caption{Utility-score scans for four cold-atom sensor scenarios. The score is
not a performance proof; it is an early rejection and ranking metric.}
\end{figure}

\section{Scientific limits}

This framework intentionally excludes laser phase noise, wavefront aberrations,
vibration rejection, magnetic systematics, atom loss, cold collisions, dead
time, navigation filtering and Allan deviation. These effects must be added in
an engineering-grade simulator.

The correct interpretation of a high score is therefore:
\begin{center}
\fbox{the phase-space budget does not already reject the design.}
\end{center}
It is not a claim that the final sensor will meet its mission specification.

\section{Industrial direction}

The credible product is software, not a speculative engine:
\begin{center}
\fbox{\texttt{PhaseSpace CAD for Quantum Sensors}}
\end{center}
Its minimum viable version should provide atom databases, phase-space density
calculation, ballistic cloud expansion, aperture survival, shot-noise
estimates, scenario ranking and CSV/JSON export. The next version should add
laser noise, vibration noise, gradiometer geometry and Allan-deviation models.

\section{Conclusion}

The same $h^3$ phase-space cell that gives little leverage for hot chemical
gases becomes useful when the industrial object is a cold quantum source. The
practical conclusion is:
\begin{equation}
\text{phase-space density}+\text{velocity dispersion}
\longrightarrow
\text{early design budget for quantum sensors}.
\end{equation}
This is a concrete, testable and conservative use of the theory.

\begin{thebibliography}{9}
\bibitem{nist-sensing}
NIST, ``Quantum Sensing Explained,'' 2026.
\url{https://www.nist.gov/quantum-information-science/quantum-sensing-explained}

\bibitem{nist-gyro}
NIST, ``Atomic Gyroscopes.''
\url{https://www.nist.gov/noac/technology/time-and-frequency/atomic-gyroscopes}

\bibitem{sandia-gravity}
Sandia National Laboratories, ``Gravimetry -- Quantum.''
\url{https://www.sandia.gov/quantum/gravimetry/}

\bibitem{darpa-roqs}
DARPA, ``From fragile to field-ready: RoQS program launches first phase,'' 2025.
\url{https://www.darpa.mil/news/2025/roqs-launches-first-phase}
\end{thebibliography}

\end{document}
